\cotefile{157-2}
\scanpage{67}
\mspage{11}

\uline{Dém} Je dis qu'il suffit de prouver, pour les trois foncteur\unsure{s} $p,q,r$, et pour les foncteurs en sens inverse
\diagram{\begin{tikzcd}[ampersand replacement=\&,column sep=large]
\underline{\mathrm{Fib}\,\mathrm{sc}}_{\unsure{0|o}}\,S
 \& \underline{\mathrm{Parf}}\,S
 \& \underline{\mathrm{Cofib}}\,\underline{\mathrm{sc}}_0\,S \\
 \& \mathrm{Cat}_{/S}\arrow[ul,"\underline{\Phi}_0"']\arrow[u,"C"]\arrow[ur,"\Psi_0"'] \&
\end{tikzcd}}
, que les six composés
\[
p\underline{\Phi}_0,\underline{\Phi}_0p,\quad q\Psi_0,\Psi_0q,\quad rC,Cr
\]
sont reliés aux foncteurs identiques (dans la catégorie pertinente) par une flèche qui est dans le localisateur \deleted{\unsure{envisagé}} correspondant. Pour \unsure{ch. un|chacun} des trois couples de foncteurs \inserted{$(p,\underline{\Phi}_0),(q,\Psi_0),(r,C)$}, cela prouvera déjà qu'en passant aux localisés, \deleted{\unsure{que}} ils induisent des équivalences quasi-inverses l'une de l'autre. Mais on a mieux, à l'aide du lemme formel suivant

\uline{Lemme} Soient $(C,W),(C',W')$ \deleted{$\unsure{(C'',W'')|(C,W)}$} \unsure{deux} catégories avec localisateurs $W\subset\operatorname{Fl}C,W'\subset\operatorname{Fl}C'$, \deleted{$\unsure{W''\subset\operatorname{Fl}C''}$}. Soit
\[ p:C\unsure{\longrightarrow|\text{---}}C' \]
un foncteur compatible avec les localisateurs, et supposons qu'il existe un foncteur $\unsure{\varphi|q}:C'\unsure{\text{---}|\longrightarrow}C$ également compatible, et tel que \deleted{\unsure{l'on}} l'on ait des flèches (à direction non précisée)
\[ \varphi p\overset{u}{\text{---}}\mathrm{id}_C,\qquad p\varphi\overset{v}{\text{---}}\mathrm{id}_{C'} \]
qui soient dans $W$ resp. $W'$ ``\unsure{argument} par argument''.

\deleted{\unsure{Alors, pour tout couple} $(A,\Sigma),\ \Sigma\subset\operatorname{Fl}(A)$, \unsure{un} \unsure{troisième} et \unsure{tout} foncteur $f:A\longrightarrow C'$}
\deleted{\unsure{Soit d'autre part}}

% archive sheet number: 66
